\section{Úvod}
\par Separace plynů --- selektivní technika dělení plynných směsí --- najde využití zejména v~odvětví energetiky, chemické výroby a ochrany životního prostředí. Jako konkrétní příklady uveďme výrobu dusíku a kyslíku ze vzduchu, separaci vodíku od dalších produktů zplyňování uhlí, dělení složek zemního plynu, oddělování oxidu uhličitého z~produktů spalování, nebo třeba zachytávání korozivních či zdraví škodlivých látek v~chemických výrobách.

\par Selektivita techniky plyne v~obecné rovině z~využitého separačního mechanismu, v~konrkétní rovině pak z~rozdílnosti hodnot příslušné veličiny, na které je mechanismus založen. Pravděpodobně nejznámnější takovou veličinou je teplota bodu varu, kdy při dostatečně velkém rozdílu lze směs dělit technikou destilace. Ať už mluvíme o frakční destilaci zkapalněného vzduchu či ropy, vždy se bude jednat o energeticky náročný, a tedy i nákladný, proces. Další často využívanou veličinou je rozdílná afinita k médiu, typicky adsorbce na zeolity. K separaci pak dochází vsádkově pomocí cyklů sorbce, desorbce a regenerace; právě regenerace pak je opět často energeticky náročná.

\subsection{Membránové seperace plynů}
\par Technika membránové separace je principielně založena na rozdílné propustnosti plynů membránou, která je ovlivněna v~první řadě velikostí a tvarem molekuly, v~druhé řadě pak afinitou plynu k adsorbci na membránu. V~porovnání s~jinými technikami nejsou membránové separace tolik energeticky náročné, jsou snadno škálovatelné a poskytují vysokou selektivitu pro rozměrově dostatečně odlišné plyny; na druhou stranu, membrány jsou náchylné k zanášení a (nejen) mechanickému poškození. 

\par Z~pohledu \textbf{materiálů} rozlišujeme membrány polymerní, keramické, nanomateriálové; homogenní, hybridní a kompozitní. Polymerní membrány jsou typicky dobře zpracovatelné, flexibilní a relativně levnější, keramické membrány nabízejí vyšší tepelnou a chemickou stabilitu, kovové membrány se vykazují vysokou selektivitou; nanomateriály díky mimořádnému efektivnímu poměru povrchu a objemu posilují mechanismus afinity. Hybridní membrány, typicky MOF (\textit{Metal-Organic Frameworks}) využívají různých funkčních příměsí; kompozitní membrány vrství různé typy membrán na sebe. 

\par Z pohledu \textbf{vnitřní stuktury} mluvíme o membránách na škále od hustých až po porézní. Husté membrány se při separaci spoléhají na rozdílnou rozpustnost a difuzivitu molekul plynu, porézní membrány obsahují síť vzájemně propojených pórů --- selektivity je dosaženo rozdílnou velikostí molekul. 

\par Fyzikálně je hnací silou procesu tlakový gradient, proti kterému stojí tloušťka membrány, separované plyny pak prostupují celou plochou membrány. K separaci dochází, pokud prostupují dostatečně rozdílnou rychlostí ---poměr propustností pak označujeme jako selektivitu.

\par Jak lze takovou technologii škálovat? Přirozeně lze zvětšit plochu membrány, kde limity pro jednu membránu vycházejí spíše z~výrobního procesu: s~rostoucí plochou roste i riziko nehomogenity výroby a v~důsledku i cena za jednotkovou výměnu na defekt. Dále se nabízí použití tenčí membrány; tloušťka však bývá omezená limitem technologie výroby, nelze také snadno vyloučit vliv tloušťky na selektivitu. V~neposlední řadě se nabízí zvýšení tlakového gradientu, ani v tomto případě však není vliv na selektivitu triviální otázkou.

\subsection{}